\begin{center}
{\LARGE Torus-fibred compact complex threefolds:\par
cotangent torsion, vanishing criteria and deformation stability}\par\medskip
{\large Companion Paper II}\par\medskip
AI-produced by GPT-6 Astra (OpenAI, ChatGPT Pro mode, effort Pro),\par
17 November 2026
\end{center}
\begin{abstract}
We identify the torsion term omitted in the singular-fibre endpoint argument
of Campana, Demailly and Peternell and prove the exact connecting-map
criterion that replaces that inference. Proper nodal and compact examples
show why the endpoint vanishings alone do not suffice. We determine the
inverse-Jacobian and conductor maps, retain multiplicities in fibre
cohomology, and construct compact torus fibrations with the specified
normalization and component classes. Integral specialization and coherent
restriction calculations give the correct component counts and Picard
conditions. We prove stability of rank-two torus fibrations, compute the
deformations of the comb and two-parameter families, and distinguish the
remaining bundle-splitting question from the established deformation results.
\end{abstract}

\section*{Main statements, status, credit and proof map}
For a reduced Cartier hypersurface $S\subset Y$, let
$N^*=\OO_Y(-S)|_S$, $V=\Omega_Y^1|_S$,
$T=\operatorname{tors}\Omega_S^1$ and $Q=\Omega_S^1/T$.
The canonical kernel $K=\ker(V\to Q)$ fits into
\[
0\longrightarrow N^*\longrightarrow K\longrightarrow T\longrightarrow0,
\qquad
0\longrightarrow K\longrightarrow V\longrightarrow Q\longrightarrow0.
\]
When $H^0(N^*\otimes A)=H^0(Q\otimes A)=0$, the ambient section
space is exactly the kernel of
$\delta_A:H^0(T\otimes A)\to H^1(N^*\otimes A)$,
as proved in \cref{cdp:st:A.1.1}.
The proper nodal example in \cref{cdp:st:B1} has a nonzero ambient
section despite both endpoint vanishings.
The replacement global criteria are \cref{cdp:st:C4a,cdp:st:C4c,cdp:st:C4f}.
The rank-two stability theorem is \cref{cdp:st:E6}; its statement includes
nonreduced deformation bases.

The contradiction between the standard complex six-sphere family and
\citet[Corollary~2.3]{CDP20} is due to
\citet[Section~1.2 and Theorem~10.5]{AC2026}.
Paper I contains the marked family, periods, discrete parameters and
standard cohomology. The present paper concerns the exact failure of the
general inference and the hypotheses under which replacement conclusions
hold. The local nodal example has precisely the hypotheses printed in its
statement; its role is not enlarged to a different global theorem.

\subsection*{Status}
\begin{tabular}{@{}p{31mm}p{120mm}@{}}
\toprule
Kind & Theorems, computations, corrections, counterexamples and credited re-derivations.\\
Claim & Checked statements and complete proofs in the body and appendices; the remaining splitting question has a separate final part.\\
Checks & Independent proof readings, original-source hypothesis matches and exact program reruns; the concordance records their individual scopes.\\
Human understanding & None claimed; no human has checked the mathematics.\\
Formalization & None; no proof-assistant certification is claimed.\\
\bottomrule
\end{tabular}

\subsection*{Known and new}
\begin{tabular}{@{}p{43mm}p{69mm}p{37mm}@{}}
\toprule
Statement & Source and exact scope & Present treatment\\
\midrule
Complex six-sphere construction and the CDP contradiction &
\citet[Theorem~10.5 and Section~1.2]{AC2026}; \citet[main theorem]{Engel2026}&
Credit retained; general cotangent and restriction analysis here.\\
Integral topology and component corrections &
\citet[Theorems~7.17, 7.22; Lemmas~2.7, 10.7; Section~10.5]{AC2026}&
Known conclusions credited; distinct complete proofs retained.\\
Inverse-Jacobian interpretation and normal-crossing criterion &
\citet[Theorems~1.1, 1.3; Proposition~3.4; Corollary~3.5]{GS2014}&
Explicit maps and full re-derivation.\\
Comb and two-parameter constructions &
The exact Viaclovsky source results identified in the individual statements&
Source construction credited; subsequent invariant and deformation claims distinguished.\\
General vanishing, occurrence and stability refinements &
The per-result comparisons in Appendix~\ref{app:II:result-index}&
Exact scopes and every distinct proof retained.\\
\bottomrule
\end{tabular}

\subsection*{Proof map}
The tangent index and Picard components isolate the global cohomological
question. Local inverse-Jacobian and conductor calculations supply the
missing fibre term. The explicit quotient, root, blowup and resolution
charts establish compact occurrence with the claimed multiplicities;
the restriction maps determine the full integral and coherent invariants.
The logarithmic Kodaira--Spencer calculation and relative graph space
give deformation stability. Root-bundle direct images then give the
comb and two-parameter deformation calculations. Complete foundational
derivations, source concordance and program instructions appear in the
appendices.

\clearpage
\section{Common conventions}
\label{sec:II:conventions}
All spaces and coherent sheaves are complex analytic unless a polynomial
ring is specified. We use $B=\PP^1_t$, $(p_1,p_2,p_0)=(0,1,\infty)$,
the ordered lattice
$\Lambda=\Z\langle\widehat\gamma,\widehat u,\widehat w,\widehat\delta\rangle$,
and the marked periods
\[
\Pi_c=\begin{pmatrix}6\mu&\tau&1&0\\\beta^{\rm part}+c&\mu&0&1\end{pmatrix},
\qquad w=e^{-\pi i c},\qquad (\beta^{\rm part}+\tau)(p_0)=0.
\]
Base meridians are counterclockwise. The analytic root decks are clockwise,
$s_j(g_jz)=e^{-2\pi i/m_j}s_j(z)$, with $g_1g_2g_0=1$;
the resulting homology monodromies are $A_1,A_2,M_0$.
The primitive invariant vectors are
$\varepsilon=(1,2,-4,0)^t$ and
$\varepsilon'=(1,3,-3,0)^t$.
For a finite affine datum we write
$v_1=r_1\varepsilon+k_1\widehat\delta$ and
$v_2=r_2\varepsilon'+k_2\widehat\delta$; the coefficients
$r_1,r_2$ are also denoted $l_1,l_2$ in the peripheral presentation.
Define
\[
 D_{\rm part}=\operatorname{Im}\beta^{\rm part}
       -\frac{6(\operatorname{Im}\mu)^2}{\operatorname{Im}\tau},
 \qquad M=\max_{B\setminus\{p_0\}}D_{\rm part}
       =-\frac3\pi\log6.
\]
The admissible parameter inequality is $\operatorname{Im}c<-M$;
the maximum of $D_{\rm part}$ is attained only at $t=1$, as proved in Paper~I,
Theorem~5.5.
The standard multiplicities are $m_1=3,m_2=4$,
$f^*p_j=m_jS_j$, and $W=f^{-1}(p_0)$ is reduced, with normalization
$\eta:\mathrm dP_6\to W$.
For general finite affine twists we use
$p=12l_0-4l_1-3l_2$, retaining $\pi_1(X)=\Z/|p|$.
The complete marked transport and construction are those of Paper I.
An unqualified fibre is its full Cartier scheme; its reduction is always
specified explicitly. Homology and cohomology retain their displayed
coefficient groups, and complex manifolds carry their complex orientations.
